% !TEX root = ../../main_power.tex
% ---------------------------------------------------------------------
%  Navigation: "Wenn die Aufgabe lautet ..." -- von der Aufgabenstellung
%  direkt zur Abschnittsnummer. Steht vor dem Inhaltsverzeichnis.
% ---------------------------------------------------------------------

\begin{tcolorbox}[zfbox, colback=black!2, colframe=black!45,
  coltitle=white, colbacktitle=black!45, title={Wenn die Aufgabe lautet \dots}]
\begin{tabularx}{\linewidth}{@{}Xr@{}}
Leistung oder Energie rechnen & \ref{sec:ref} \\
Batterielaufzeit abschätzen & \ref{sec:ref} \\
Verbrauch senken: welche Schraube zuerst & \ref{form:ptotal} \\
Statisch gegen dynamisch unterscheiden & \ref{sec:dynamisch} \\
Warum steigt der Verbrauch beim Schrumpfen & \ref{sec:dynamisch} \\
Power Gating oder Clock Gating & \ref{rez:gating} \\
Hardware-Massnahmen aufzählen & \ref{sec:pwrhw} \\
Software-Massnahmen aufzählen & \ref{sec:pwrsw} \\
Welcher Kinetis-Modus zu welcher Weckquelle & \ref{rez:modi} \\
Gear Shifting implementieren & \ref{rez:gearshift} \\
FreeRTOS schlafen legen & \ref{rez:idlehook} \\
Warum reicht der Idle Hook nicht & \ref{rez:tickless} \\
Messaufbau und Messfehler & \ref{sec:messen} \\
Warum ist die Messung unplausibel & \ref{falle:pwr} \\
\end{tabularx}
\vspace{2pt}
\textcolor{accent}{\rule{7pt}{7pt}}\,Formel \;
\textcolor{red!65}{\rule{7pt}{7pt}}\,Falle \;
\textcolor{green!40!black}{\rule{7pt}{7pt}}\,Rezept \;
\textcolor{black!55}{\rule{7pt}{7pt}}\,Beispiel \;
-- der letzte Abschnitt ist das Stichwortregister.
\end{tcolorbox}
